% Part: normal-modal-logic
% Chapter: filtrations
% Section: introduction

\documentclass[../../../include/open-logic-section]{subfiles}

\begin{document}

\olfileid{nml}{fil}{int}

\olsection{سريزه}

د يوه منطق په اړه يوه مهمه پوښتنه دا وي چې آيا هغه د پرېکړې وړ دے، يعنې آيا
داسې اغېزمنه کړنلاره شته چې د «آيا دا !!{formula} اعتبار لري؟» پوښتنې ته ځواب
ورکړي۔ بياني منطق د پرېکړې وړ دے: د رښتياوالي د جدول په جوړولو په اغېزمن ډول
ازمېيلے شو چې آيا !!a{formula} تاوتولوژي ده؛ او د ټاکلي !!{formula} لپاره
جدول متناهي دے۔ خو د مودالي فارمول د ټولو مدلونو په نسبت رښتياوالے په ښکاره
ډول نۀ شو ازمېيلے، ځکه مدلونه نامتناهي ډېر دي۔ د يوه ټاکلي فارمول لپاره
اړوند ټول متناهي مدلونه په لړ کښې راوستلے شو، ځکه يوازې د هغو
!!{propositional variable}s لپاره د نړۍو د فرعي سټونو ټاکنه اړونده ده چې
رښتيا هم په !!{formula} کښې راځي۔ که د لاسرسي اړيکه ثابته وي، د $V(p)$
بېلابېلې ممکنې ټاکنې د~$W$ ټول فرعي سټونه دي؛ او که $\card{W} = n$ وي،
نو د هغو شمېر $2^n$ دے۔ که زموږ !!{formula}~$!A$، $m$ !!{propositional
  variable}s ولري، نو د~$n$ نړۍو لپاره $2^{nm}$ بېلابېل مدلونه شته۔ د هر
مدل لپاره دا ازمېيلے شو چې آيا~$!A$ په ټولو نړۍو کښې رښتينے دے، يوازې په هره
نړۍ کښې د~$!A$ د رښتياوالي قيمت په محاسبه کولو سره۔ البته د لاسرسي ټولې ممکنې
اړيکې هم بايد وګورو، خو پر~$n$ نړۍو اړيکې هم متناهي ډېرې دي (په ځانګړي ډول
د $W \times W$ د فرعي سټونو شمېر، يعنې $2^{n^2}$)۔

که زموږ پام منطق~\Log{K} ته نۀ وي، بلکې داسې منطق ته وي چې د مدلونو د يوه
ټولګي په وسيله تعريف شوے وي (د بېلګې په توګه، انعکاسي متعدي مدلونه)، نو بايد
دا هم ازمېيلے شو چې د لاسرسي اړيکه د غوښتل شوي ډول ده که نه۔ کله چې زموږ
چوکاټونه د مودالي !!{formula}s له لارې د تعريف وړ وي، دا ازموينه کولے شو
(د بېلګې په توګه، په چوکاټ کښې د \Ax{T} او \Ax{4} د اعتبار په ازمېيلو)۔ نو
لار دا ښکاري چې ټول متناهي چوکاټونه په لړ کښې واخلو، هر يو وازمېيو چې آيا
د غوښتل شوي ټولګي چوکاټ دے، بيا پر هغه چوکاټ ټول ممکن مدلونه راوړو او
وازمېيو چې آيا $!A$ په هر يوه کښې رښتينے دے۔ که داسې نۀ وي، ودرېږو: $!A$
په غوښتل شوي مدل‌ټولګي کښې اعتبار نۀ لري۔

په دې لاره کښې ستونزه شته: نۀ پوهېږو چې لټون کله، او که هېڅکله، درولے شو۔
که !!{formula} متناهي ضدمدل ولري، کړنلاره به ئې ومومي۔ خو که متناهي ضدمدل
ونۀ لري، ځواب به رانۀ کړي۔ کېدے شي !!{formula} معتبر وي (يعنې هېڅ ضدمدل
ونۀ لري)، يا يوازې نامتناهي ضدمدل ولري چې موږ ئې هېڅکله نۀ ګورو۔ دا ستونزه
هله هوارېدے شي چې وښيو هر هغه !!{formula} چې ضدمدل لري، متناهي ضدمدل هم
لري۔ په دې صورت کښې وايو چې منطق د \emph{متناهي مدل خاصيت} لري۔

خو څنګه وښيو چې يو منطق د متناهي مدل خاصيت لري؟ يوه لار دا ده چې د~$!A$
يو نامتناهي (ضد)مدل په متناهي مدل بدل کړو۔ که دا کار وشي، نو هر کله چې داسې
مدل وي چې $!A$ پکښې رښتينے نۀ وي، له هغه نه جوړ شوے متناهي مدل به هم $!A$
نارښتينے کړي۔ دا متناهي مدل به زموږ د ټولو متناهي مدلونو په لړ کښې راشي؛ نو
د هر نامعتبر فارمول په اړه به بالاخره معلومه کړو چې نامعتبر دے۔ خو که
فارمول معتبر وي، دا لټونيزه کړنلاره به نۀ درېږي۔ که سربېره پر دې وښيو چې
زموږ د کړنلارې له لارې ترلاسه کېدونکے متناهي مدل يوه اعظمي کچه لري، او دا کچه
يوازې په !!{formula}~$!A$ پورې تړلې ده، نو د ضدمدل د لټون لپاره به پوه شو
چې تر کومې کچې متناهي مدلونه وازمېيو۔ که تر هغه وخته ضدمدل ونۀ موندل شي،
هېڅ ضدمدل نشته۔ بيا به زموږ کړنلاره د هر !!{formula}~$!A$ لپاره د «آيا
$!A$ اعتبار لري؟» پوښتنه په رښتيا پرېکړه کړي۔

د نامتناهي جوړښتونو په متناهي جوړښتونو د بدلولو يوه ګټوره تګلاره دا ده چې
د جوړښت هغه !!{element}s «يو شان» وګڼو چې له اړوندو پلوو يو شان چلند کوي۔
که په~$\mModel{M}$ کښې نامتناهي ډېرې نړۍ له اړوندو پلوو يو شان چلند کوي،
کېدے شي \emph{ټولې} نړۍ د دغو نړۍو په متناهي شمېر ټولګيو کښې راټولې کړو؛
هر ټولګے په خپله متناهي يا نامتناهي کېدے شي۔ (د سرچينې په دې تشريحي جمله
کښې د هر ټولګي نامتناهي‌والے لازم بلل شوے؛ د متناهي نوي مدل لپاره يوازې
د ټولګيو متناهي شمېر لازم دے، لکه چې ورپسې رسمي تعريف ښيي۔) په بله وينا،
د نړۍو سټ بايد په مناسب ډول ووېشو، داسې چې ټولې نړۍ پکښې راشي او د وېش
ټولګي يوازې متناهي ډېر وي۔ بيا يو نوے مدل~$\mModel{M^*}$ تعريفوو چې نړۍې
ئې همدا ټولګي وي۔ د پخواني مدل متناهي ډېر ټولګي په نوي مدل کښې متناهي ډېرې
نړۍې راکوي، يعنې متناهي مدل۔ د هغې وېش‌برخې نوم چې نړۍ~$w$ پکښې ده، $[w]$
کېږدو۔ غواړو $\mSat{M}{!A}[w]$ هله او يوازې هله وي چې
$\mSat{M^*}{!A}[{[w]}]$ وي، ځکه که پخوانے مدل د~$!A$ ضدمدل وي، نو
نوے مدل هم بايد ضدمدل پاتې شي۔ د دې لپاره وېش او د~$\mModel{M^*}$
د لاسرسي اړيکه په سمه توګه تعريفول پکار دي۔

د دې د څرنګوالي د ليدو لپاره، لومړے داسې ساده حالت وګڼو چې هره نړۍ هرې
نړۍ ته لاسرسی لري۔ په دې حالت کښې $\mSat{M}{\Box !B}[w]$ هله او يوازې
هله چې د هر $v \in W$ لپاره $\mSat{M}{!B}[v]$؛ په~$\mModel{M^*}$
کښې همداسې ده، خو د $[w]$ او $[v]$ له لارې۔ (سرچينه دلته «د لاسرسي
اړيکه نشته» او د بکس د شرط لپاره «يوه نړۍ» وايي؛ د ورپسې «هرې نړۍ»
استقرايي استدلال لپاره عمومي لاسرسی او د هرې نړۍ شرط لازم دي۔) د لومړۍ
ازموينې په توګه ووايو چې دوه نړۍې $u$ او $v$ هله همارزې دي (په يوه
وېش‌برخه کښې دي) چې په~$\mModel{M}$ کښې د ټولو !!{propositional variable}s
په اړه يوه خوله وي، يعنې $\mSat{M}{p}[u]$ هله او يوازې هله چې
$\mSat{M}{p}[v]$۔ پرېږدئ چې $V^*(p) = \Setabs{[w]}{\mSat{M}{p}[w]}$۔
غواړو وښيو چې $\mSat{M}{!A}[w]$ هله او يوازې هله چې
$\mSat{M^*}{!A}[{[w]}]$۔ دا په استقرا ثابتوو: بنسټيز صورت $!A
\equiv p$ دے۔ لومړے فرض کړو چې $\mSat{M}{p}[w]$۔ نو د تعريف له مخې
$[w] \in V^*(p)$، او ځکه $\mSat{M^*}{p}[{[w]}]$۔ (د سرچينې په همدې
بنسټيز ګام کښې د ارزښت ټاکنې د متغير نښه لوېدلې؛ دلته د هماغه متغير
ارزښت‌سټ په څرګند ډول ليکل شوے۔) اوس فرض کړو چې
$\mSat{M^*}{p}[{[w]}]$۔ دا معنا لري چې $[w] \in V^*(p)$، يعنې د~$w$
سره همارزې يوې نړۍ~$v$ لپاره $\mSat{M}{p}[v]$۔ خو «$w$ له $v$
سره همارز دے» دا معنا لري چې «$w$ او $v$ ټول هماغه !!{propositional variable}s
رښتيني کوي»؛ نو $\mSat{M}{p}[w]$۔ اوس استقرايي ګام، د بېلګې په توګه،
$!A \ident \lnot !B$۔ بيا $\mSat{M}{\lnot !B}[w]$ هله او يوازې هله چې
$\mSat/{M}{!B}[w]$، هله او يوازې هله چې $\mSat/{M^*}{!B}[{[w]}]$
(د استقرا فرضيې له مخې)، هله او يوازې هله چې
$\mSat{M^*}{\lnot !B}[{[w]}]$۔ نور نامودالي عملګر هم همداسې دي۔ د
$\Box$ لپاره هم کار کوي: فرض کړو $\mSat{M^*}{\Box
  !B}[{[w]}]$۔ دا معنا لري چې د هر $[u]$ لپاره
$\mSat{M^*}{!B}[{[u]}]$۔ د استقرا فرضيې له مخې، د هر $u$ لپاره
$\mSat{M}{!B}[u]$۔ نو $\mSat{M}{\Box !B}[w]$۔

په عمومي صورت کښې، چې د~$\mModel{M^*}$ د لاسرسي اړيکه هم بايد تعريف
کړو، چارې پېچلې دي۔ يو مدل~$\mModel{M^*}$ ته \emph{فلټريشن} وايو که
اړيکه~$R^*$ ئې هغه شرطونه پوره کړي چې پورته استقرايي ثبوت ته اړ دي۔
بيا هر فلټريشن~$\mModel{M^*}$، $!A$ په $[w]$ کښې هله او يوازې هله
رښتينے کوي چې $\mModel{M}$، $!A$ په~$w$ کښې رښتينے کوي۔ خو اوس بايد دا
هم وښيو چې فلټريشنونه \emph{شته}، يعنې~$R^*$ داسې تعريفولے شو چې اړين
شرطونه پوره کړي۔ د دې کار لپاره دا هم لازم دي چې نړۍې $u$ او $v$ يوازې
د ټولو !!{propositional variable}s په اړه د يو شان رښتياوالي له امله
همارزې ونۀ ګڼو؛ بلکې د~$!A$ د ټولو فرعي !!{formula}s په اړه ئې هم
رښتياوالے يو شان وي۔ ځکه چې $!A$ يوازې متناهي ډېر فرعي !!{formula}s
لري، دا بيا هم د فلټريشن متناهي‌والے تضمينوي۔ د فلټريشن د تعريف يوازې
يوه لار نشته؛ او د دې لپاره چې د فلټريشن د لاسرسي اړيکه غوښتل شوي
خاصيتونه (لکه انعکاس، تعدّي او نور) ولري، د~$R^*$ په تعريف کښې
له فکره کار اخيستل پکار دي۔

\end{document}
